\chapter{电势与电容}

电场力和万有引力一样，都是\emph{保守力}。这意味着：电荷从一点移动到另一点时，电场力做功只由起点和终点决定，而与路径的具体形状无关。于是我们可以像力学中引入重力势能那样，在静电学中引入\textbf{电势}这一标量函数。它把原本带有方向的电场信息，压缩为一个更容易计算、也更容易与能量观点结合的量。

本章有三条主线：
\begin{enumerate}
  \item 电势差来自电场沿路径的线积分；
  \item 电场是电势在空间中的变化率，即 $\vb E=-\grad V$；
  \item 电容与电场能量都能由电势观点自然推出。
\end{enumerate}

从高中物理的角度看，电势常被理解为“单位正电荷的电势能”；从微积分的角度看，它更像是一个\emph{势函数}：知道空间中每一点的 $V(\vb r)$，就能用求导恢复电场，用积分计算总能量。

\section{电势}

\subsection{电势差的线积分定义}

设试探电荷 $q_0$ 从点 $A$ 沿某条曲线移动到点 $B$，路径元记为 $\dd\vb l$。电场力的元功为
\[
\dd W = q_0\,\vb E\cdot \dd\vb l,
\]
因此总功为
\[
W_{A\to B}=q_0\int_A^B \vb E\cdot \dd\vb l.
\]

若只考虑静电场，由于它是保守场，电场力做功等于电势能的减少：
\[
W_{A\to B}=-(U_B-U_A).
\]
两边同除以 $q_0$，便得到单位正电荷的势能变化，也即电势差。

\begin{definition}
点 $A,B$ 之间的\textbf{电势差}定义为
\[
V_B-V_A=-\int_A^B \vb E\cdot \dd\vb l.
\]
若选定某个参考点 $O$ 使 $V(O)=0$，则点 $P$ 的\textbf{电势}为
\[
V(P)-V(O)=-\int_O^P \vb E\cdot \dd\vb l.
\]
在静电学中，常取无穷远处为零势点，即 $V(\infty)=0$。
\end{definition}

这个定义强调了两件事：
\begin{itemize}
  \item 电势是\textbf{标量}，没有方向；
  \item 电势差的定义本质上是\textbf{线积分}，而静电场的特殊性在于该积分与路径无关。
\end{itemize}

\begin{theorem}
在静电场中，若空间区域内不存在随时间变化的磁场，则对任意闭合回路 $\Gamma$ 都有
\[
\oint_\Gamma \vb E\cdot \dd\vb l=0.
\]
因此从一点到另一点的线积分只与端点有关，电势函数 $V(\vb r)$ 可以在该区域内定义。
\end{theorem}

\begin{proof}
静电场满足 $\curl \vb E=\vb 0$。由斯托克斯公式，
\[
\oint_\Gamma \vb E\cdot \dd\vb l
=\iint_S (\curl \vb E)\cdot \dd\vb S=0.
\]
所以任意两条连接同一对端点的路径，其线积分之差为零，故线积分只由端点决定。这正是势函数存在的充分条件。
\end{proof}

\begin{remark}
“电势高低”与“电势能大小”并不完全相同。对正电荷，电势越高，势能越高；对负电荷，情形恰好相反，因为 $U=qV$ 中的 $q$ 可以为负。
\end{remark}

\subsection{电势的物理意义}

若把单位正电荷从参考点缓慢移到空间某点 $P$，外力所做的功（准静态条件下）就等于该点的电势。于是对于任意电荷 $q$，其电势能为
\[
U=qV.
\]
这说明：电势本身是“场的性质”，而势能是“电荷与场共同决定的量”。

\begin{example}
在匀强电场中，电场方向沿 $x$ 轴正方向，$\vb E=E_0\uvec{x}$。求两点 $A(x_A)$、$B(x_B)$ 之间的电势差。
\end{example}

\begin{proof}[解]
取路径沿 $x$ 轴，则 $\dd\vb l=\dd x\,\uvec{x}$，从而
\[
V_B-V_A=-\int_{x_A}^{x_B} E_0\,\dd x=-E_0(x_B-x_A).
\]
所以电势沿电场方向线性降低。若 $x_B>x_A$，则 $V_B<V_A$。
\end{proof}

这一结论可以口头记成：\textbf{顺着电场方向走，电势降低；逆着电场方向走，电势升高。}

\section{电势与电场的关系}

电势是全局量，电场是局部量。二者之间的联系由微积分中的导数给出。

\subsection{一维情形：导数形式}

若电势只随 $x$ 变化，由定义有
\[
V(x+\dd x)-V(x)=-E_x\,\dd x.
\]
两边除以 $\dd x$ 并取极限，得
\[
E_x=-\dv{V}{x}.
\]
这说明电场强度等于电势对位置变化率的负值。

\subsection{三维情形：梯度形式}

在三维空间中，
\[
\dd V=-\vb E\cdot \dd\vb l.
\]
另一方面，函数 $V(x,y,z)$ 的全微分为
\[
\dd V=\frac{\partial V}{\partial x}\dd x+\frac{\partial V}{\partial y}\dd y+\frac{\partial V}{\partial z}\dd z=\grad V\cdot \dd\vb l.
\]
对任意位移 $\dd\vb l$ 比较两式，得到静电学中极其重要的关系。

\begin{law}
静电场与电势的关系为
\[
\vb E=-\grad V.
\]
在直角坐标系中即
\[
E_x=-\pdv{V}{x},\qquad E_y=-\pdv{V}{y},\qquad E_z=-\pdv{V}{z}.
\]
\end{law}

该公式有非常直观的几何含义：电场方向指向电势下降最快的方向，而电场大小等于单位长度上的最大降势率。

\subsection{等势面}

\begin{definition}
若空间中某个曲面上所有点的电势都相等，则该曲面称为\textbf{等势面}；二维截面中的对应曲线称为\textbf{等势线}。
\end{definition}

沿等势面移动时 $\dd V=0$，于是
\[
\vb E\cdot \dd\vb l=0.
\]
因此电场总与等势面垂直。这与地形中的等高线完全类似：坡度最大的方向垂直于等高线，电势下降最快的方向则是电场方向。

\begin{remark}
等势面不能相交。否则同一点将被赋予两个不同的电势值，这与电势作为标量场的定义矛盾。
\end{remark}

\begin{figure}[H]
\centering
\begin{tikzpicture}[scale=0.95]
  \draw[->, thick] (-3.5,0) -- (3.8,0) node[right] {$x$};
  \draw[->, thick] (0,-3.5) -- (0,3.8) node[above] {$y$};
  \fill[red!70] (0,0) circle (0.08) node[below right] {$+q$};
  \foreach \r/\lab in {0.8/$V_1$,1.5/$V_2$,2.2/$V_3$,2.9/$V_4$}{
    \draw[blue!70] (0,0) circle (\r);
    \node[blue!70] at ({0.72*\r},{0.72*\r}) {\small \lab};
  }
  \foreach \ang in {0,30,...,330}{
    \draw[orange!85,->] (0,0) -- (\ang:3.15);
  }
  \node at (0,-4.1) {蓝色为等势线，橙色为电场线，二者处处垂直};
\end{tikzpicture}
\caption{点电荷周围的等势线与电场线}
\end{figure}

\begin{example}
已知一维电势分布为
\[
V(x)=V_0-ax+bx^2,
\]
其中 $a,b$ 为正常数。求电场 $E_x(x)$，并判断 $x=\dfrac{a}{2b}$ 两侧电场方向。
\end{example}

\begin{proof}[解]
由 $E_x=-\dv*{V}{x}$，得
\[
E_x=a-2bx.
\]
当 $x<\dfrac{a}{2b}$ 时，$E_x>0$，电场沿 $x$ 正方向；当 $x>\dfrac{a}{2b}$ 时，$E_x<0$，电场沿 $x$ 负方向；在 $x=\dfrac{a}{2b}$ 处电场为零。也就是说，该点是电势函数的驻点，对应电场方向发生改变。
\end{proof}

\section{点电荷与电荷分布的电势}

\subsection{点电荷的电势}

对位于原点的点电荷 $q$，其电场为
\[
\vb E=\frac{1}{4\pi\varepsilon_0}\frac{q}{r^2}\uvec{r}.
\]
取无穷远处为零势点，沿径向从 $\infty$ 积分到 $r$，则
\[
V(r)=-\int_\infty^r \vb E\cdot \dd\vb l
=-\int_\infty^r \frac{1}{4\pi\varepsilon_0}\frac{q}{r'^2}\dd r'.
\]
计算得
\[
V(r)=\frac{1}{4\pi\varepsilon_0}\frac{q}{r}.
\]

\begin{law}
点电荷 $q$ 在距离其 $r$ 处产生的电势为
\[
V(r)=\frac{1}{4\pi\varepsilon_0}\frac{q}{r},
\]
并满足叠加原理：若有多个点电荷，则总电势等于各点电势的代数和。
\end{law}

于是离散电荷系的电势写为
\[
V(\vb r)=\frac{1}{4\pi\varepsilon_0}\sum_i \frac{q_i}{\abs{\vb r-\vb r_i}}.
\]
注意这是\textbf{代数和}而不是矢量和，因此通常比电场更容易计算。

\begin{figure}[htbp]
\centering
\begin{tikzpicture}
\begin{axis}[
  width=0.75\textwidth, height=0.5\textwidth,
  xlabel={距离 $r$}, ylabel={电势 $V$},
  axis lines=middle,
  xmin=0, xmax=4.2, ymin=0, ymax=2.4,
  grid=major, grid style={gray!30},
  thick,
]
\addplot[blue, domain=0.6:4, samples=180] {1/x};
\end{axis}
\end{tikzpicture}
\caption{点电荷电势随距离按 $1/r$ 下降}
\end{figure}

\begin{example}
两个点电荷 $+q$ 与 $-q$ 分别放在 $x=\pm a$ 处。求原点的电势，并判断原点电场的方向。
\end{example}

\begin{proof}[解]
原点到两电荷的距离都为 $a$，故电势为
\[
V(0)=\frac{1}{4\pi\varepsilon_0}\qty(\frac{q}{a}-\frac{q}{a})=0.
\]
但电场并不为零。原点处由 $+q$ 产生的电场指向远离正电荷的方向，即向左；由 $-q$ 产生的电场指向负电荷本身，也向左。二者叠加后，原点电场沿 $x$ 负方向。

因此，\textbf{电势为零并不意味着电场为零}；电势取决于标量叠加，电场取决于矢量叠加。
\end{proof}

\subsection{连续分布电荷的电势积分}

若电荷连续分布，则把电荷元 $\dd q$ 看成无穷多个点电荷，利用叠加原理便得
\[
\dd V=\frac{1}{4\pi\varepsilon_0}\frac{\dd q}{R},
\]
其中 $R$ 是场点到电荷元的距离。积分后有
\[
V=\frac{1}{4\pi\varepsilon_0}\int \frac{\dd q}{R}.
\]

根据分布类型不同，可分别写成
\begin{align*}
V(\vb r)&=\frac{1}{4\pi\varepsilon_0}\int \frac{\lambda\,\dd l'}{\abs{\vb r-\vb r'}},\\
V(\vb r)&=\frac{1}{4\pi\varepsilon_0}\iint \frac{\sigma\,\dd S'}{\abs{\vb r-\vb r'}},\\
V(\vb r)&=\frac{1}{4\pi\varepsilon_0}\iiint \frac{\rho\,\dd \tau'}{\abs{\vb r-\vb r'}}.
\end{align*}

这些公式看上去像“库仑定律的积分版”，但它们有一个重要优势：分母中只有距离 $R$，没有方向分量，所以很多对称问题先算电势再求电场会更简洁。

\begin{example}
半径为 $R$ 的均匀带电圆环，总电荷为 $Q$。求圆环轴线上距离圆心 $x$ 处的电势，并由电势求轴向电场。
\end{example}

\begin{proof}[解]
圆环上任意电荷元到轴线上场点的距离都相同，均为
\[
\sqrt{x^2+R^2}.
\]
因此
\[
V(x)=\frac{1}{4\pi\varepsilon_0}\int \frac{\dd q}{\sqrt{x^2+R^2}}
=\frac{1}{4\pi\varepsilon_0}\frac{Q}{\sqrt{x^2+R^2}}.
\]
再由
\[
E_x=-\dv{V}{x}
\]
得到
\[
E_x=\frac{1}{4\pi\varepsilon_0}\frac{Qx}{(x^2+R^2)^{3/2}}.
\]
这里清楚地体现出“先积分求势，再求导得场”的微积分思路。
\end{proof}

作为另一类常见积分，长度为 $2L$ 的均匀带电细棒沿 $x$ 轴放置、线电荷密度为 $\lambda$ 时，其垂直平分线上距中点 $a$ 处的电势可写成
\[
V(a)=\frac{1}{4\pi\varepsilon_0}\int_{-L}^{L}\frac{\lambda\,\dd x}{\sqrt{x^2+a^2}}
=\frac{\lambda}{2\pi\varepsilon_0}\ln\qty(\frac{L+\sqrt{L^2+a^2}}{a}).
\]
当 $L\to \infty$ 时，上式发散，说明无限长带电直线若取无穷远为零势点会遇到参考点选择问题，这也是“绝对电势依赖参考零点”的一个例子。

\subsection{从电势反求电场}

一旦求出 $V(\vb r)$，就可以通过
\[
\vb E=-\grad V
\]
恢复电场。这常比直接对电场做分量积分更高效。对称性越强、积分变量越多时，这种“先势后场”的策略就越有威力。

\section{电容器}

\subsection{电容的定义}

如果一个导体或一组导体能够在一定电势差下储存电荷，我们就说它具有电容。最常见的是两导体构成的\textbf{电容器}。

\begin{definition}
对一对带等量异种电荷 $\pm Q$、电势差为 $U$ 的导体系统，其\textbf{电容}定义为
\[
C=\frac{Q}{U}.
\]
在电容器问题中，$U$ 常写作两极板间电压，因此也常见写法 $C=Q/V$。
\end{definition}

电容刻画的是“单位电压下能存多少电荷”，只由几何形状、相对位置以及介质性质决定，与实际带了多少电荷无关。

\subsection{平行板电容器}

设两块面积为 $S$、间距为 $d$ 的平行导体板带电 $\pm Q$，忽略边缘效应。面电荷密度为
\[
\sigma=\frac{Q}{S}.
\]
两板间电场近似为匀强电场：
\[
E=\frac{\sigma}{\varepsilon_0}=\frac{Q}{\varepsilon_0 S}.
\]
故两板间电势差为
\[
U=Ed=\frac{Qd}{\varepsilon_0 S}.
\]
代入电容定义，便得
\[
C=\frac{Q}{U}=\frac{\varepsilon_0 S}{d}.
\]

\begin{law}
真空中平行板电容器的电容为
\[
C=\frac{\varepsilon_0 S}{d},
\]
其中 $S$ 为极板正对面积，$d$ 为板间距离。增大面积会增大电容，增大间距会减小电容。
\end{law}

\begin{figure}[H]
\centering
\begin{tikzpicture}[scale=1.0]
  \fill[gray!35] (-3,2) rectangle (-2.6,-2);
  \fill[gray!35] (2.6,2) rectangle (3,-2);
  \node at (-2.8,2.35) {$+Q$};
  \node at (2.8,2.35) {$-Q$};
  \foreach \y in {-1.5,-0.75,0,0.75,1.5}{
    \draw[blue!75,->,thick] (-2.45,\y) -- (2.45,\y);
  }
  \draw[<->] (-2.6,-2.5) -- (2.6,-2.5) node[midway, fill=white] {$d$};
  \draw[decorate,decoration={brace,mirror,amplitude=6pt}] (-3,-2.15) -- (-2.6,-2.15) node[midway,yshift=-0.55cm] {$\text{极板}$};
  \node at (0,2.35) {$\vb E$ 近似均匀};
\end{tikzpicture}
\caption{平行板电容器及其近似匀强电场}
\end{figure}

\begin{remark}
平行板电容器公式建立在“板间距离远小于板的线度、可忽略边缘效应”的前提下。若极板很小或间距很大，电场会向外弯曲，公式只剩近似意义。
\end{remark}

\begin{example}
一个空气平行板电容器的面积增为原来的 $2$ 倍、板间距离减为原来的 $\dfrac{1}{2}$。求其电容变为原来的多少倍。
\end{example}

\begin{proof}[解]
由 $C=\varepsilon_0 S/d$，得
\[
\frac{C'}{C}=\frac{(2S)/(d/2)}{S/d}=4.
\]
所以新电容是原来的 $4$ 倍。这表明电容本质上是一个几何量：增大有效面积、缩小板间距离都会增强储电能力。
\end{proof}

\section{电场能量}

\subsection{充电过程中外力做功}

把一个电容器从不带电逐渐充到总电荷为 $Q$。若某一时刻极板上已有电荷 $q$，其电压为
\[
U=\frac{q}{C}.
\]
这时再搬运微元电荷 $\dd q$ 到极板上，外力需做元功
\[
\dd W=U\,\dd q=\frac{q}{C}\dd q.
\]
从 $0$ 积分到 $Q$，得总能量
\[
W=\int_0^Q \frac{q}{C}\dd q=\frac{Q^2}{2C}.
\]
结合 $Q=CU$，还能写成
\[
W=\frac{1}{2}QU=\frac{1}{2}CU^2.
\]

若把电容器通过电阻接到恒定电压源上，极板电荷会按指数规律逐渐逼近终值 $Q_0$，其典型充电曲线为
\[
Q(t)=Q_0\left(1-\mathrm{e}^{-t/(RC)}\right).
\]

\begin{figure}[htbp]
\centering
\begin{tikzpicture}
\begin{axis}[
  width=0.75\textwidth, height=0.5\textwidth,
  xlabel={时间 $t/(RC)$}, ylabel={电荷 $Q/Q_0$},
  axis lines=middle,
  xmin=0, xmax=5.2, ymin=0, ymax=1.08,
  grid=major, grid style={gray!30},
  thick,
]
\addplot[blue, domain=0:5, samples=180] {1 - exp(-x)};
\end{axis}
\end{tikzpicture}
\caption{电容器充电时极板电荷随时间逐渐接近饱和值}
\end{figure}

\begin{law}
电容器储存的电场能量为
\[
W=\frac{Q^2}{2C}=\frac{1}{2}QU=\frac{1}{2}CU^2.
\]
\end{law}

\begin{figure}[htbp]
\centering
\begin{tikzpicture}
\begin{axis}[
  width=0.75\textwidth, height=0.5\textwidth,
  xlabel={电压 $U$}, ylabel={储能 $W$},
  axis lines=middle,
  xmin=0, xmax=4.2, ymin=0, ymax=8.8,
  grid=major, grid style={gray!30},
  thick,
]
\addplot[blue, domain=0:4, samples=180] {0.5*x^2};
\end{axis}
\end{tikzpicture}
\caption{电容器储能随电压呈抛物线增长}
\end{figure}

\subsection{能量密度}

以真空平行板电容器为例，板间体积为 $Sd$，而
\[
W=\frac{1}{2}CU^2=\frac{1}{2}\cdot \frac{\varepsilon_0 S}{d}\cdot (Ed)^2
=\frac{1}{2}\varepsilon_0 E^2\,Sd.
\]
于是单位体积中的能量为
\[
u=\frac{W}{Sd}=\frac{1}{2}\varepsilon_0 E^2.
\]
这就得到真空中电场能量密度的表达式。

\begin{law}
真空电场的能量密度为
\[
u=\frac{1}{2}\varepsilon_0 E^2.
\]
若电场分布不均匀，则总能量由体积分给出：
\[
W=\iiint \frac{1}{2}\varepsilon_0 E^2\,\dd \tau.
\]
\end{law}

这意味着电场能量并不是“神秘地附着在电荷上”，而是分布在整个有电场的空间中。场强大的地方，能量密度也大。

对接在恒定电压 $U$ 的真空平行板电容器，板间电场近似为
\[
E=\frac{U}{d},
\]
故其能量密度为
\[
u=\frac{1}{2}\varepsilon_0 E^2=\frac{1}{2}\varepsilon_0\frac{U^2}{d^2}.
\]
若电压保持不变而减小板距，则电场变强，能量密度增大。同时电容 $C=\varepsilon_0 S/d$ 也增大，总能量
\[
W=\frac{1}{2}CU^2=\frac{1}{2}\frac{\varepsilon_0 S}{d}U^2
\]
也随 $d$ 减小而增大。

\subsection{一般电荷分布的能量表达式}

对任意静电系统，把电荷从无穷远逐步搬运进来，可以得到总电场能量还可写为
\[
W=\frac{1}{2}\int V\,\dd q.
\]
若为连续体分布，则写作
\[
W=\frac{1}{2}\iiint \rho V\,\dd \tau.
\]
这里的 $\dfrac{1}{2}$ 用来避免把任意一对电荷相互作用能计算两次。这一公式与 $\iiint \dfrac{1}{2}\varepsilon_0 E^2\dd\tau$ 本质等价，只是出发点不同：一个从“搬运电荷”看，一个从“空间储能”看。

\section{电介质}

\subsection{极化现象}

在真空中，电场线可以看作直接穿过空间；但若在电场中放入绝缘材料，材料内部的正负电荷会发生微小相对位移，形成许多微观电偶极子，这种现象称为\textbf{极化}。

\begin{definition}
电介质在外电场作用下，内部束缚正负电荷发生微小分离，从而形成宏观偶极化的现象，称为\textbf{极化}；能够发生这种极化而又不形成持续导电电流的材料称为\textbf{电介质}。
\end{definition}

极化后，介质表面会出现束缚电荷，它们产生附加电场，方向通常与原外电场相反，于是材料内部的合场强减小。

\subsection{介电常数}

实验表明，对许多线性、各向同性电介质，插入介质后电场变为原来的 $1/\kappa$，其中 $\kappa$ 称为相对介电常数。于是可把真空介电常数 $\varepsilon_0$ 替换为
\[
\varepsilon=\kappa\varepsilon_0.
\]
因此许多公式都可形式不变地推广：
\[
C=\frac{\varepsilon S}{d}=\kappa\frac{\varepsilon_0 S}{d}=\kappa C_0.
\]
这里 $C_0$ 为真空中的电容。

\begin{law}
在极板间完全充满均匀电介质的平行板电容器中，电容变为
\[
C=\kappa \frac{\varepsilon_0 S}{d}=\kappa C_0.
\]
介电常数越大，电容越大。
\end{law}

\begin{figure}[H]
\centering
\begin{tikzpicture}[scale=0.98]
  \begin{scope}[xshift=-4.8cm]
    \fill[gray!35] (-1.2,1.8) rectangle (-0.8,-1.8);
    \fill[gray!35] (0.8,1.8) rectangle (1.2,-1.8);
    \foreach \y in {-1.2,-0.4,0.4,1.2}{
      \draw[blue!75,->,thick] (-0.7,\y) -- (0.7,\y);
    }
    \node at (0,-2.35) {真空};
  \end{scope}
  \begin{scope}[xshift=4.8cm]
    \fill[gray!35] (-1.2,1.8) rectangle (-0.8,-1.8);
    \fill[gray!35] (0.8,1.8) rectangle (1.2,-1.8);
    \fill[green!20] (-0.8,1.8) rectangle (0.8,-1.8);
    \foreach \y in {-1.2,-0.4,0.4,1.2}{
      \draw[blue!75,->,thick] (-0.7,\y) -- (0.7,\y);
      \draw[orange!70,->,thick] (0.45,\y+0.12) -- (-0.45,\y+0.12);
    }
    \node at (0,-2.35) {充满电介质};
  \end{scope}
\end{tikzpicture}
\caption{电介质极化会产生反向附加电场，使合场减弱}
\end{figure}

\begin{example}
一个平行板电容器先与电源断开，再在板间完全插入相对介电常数为 $\kappa$ 的电介质。求电压、场强和储能如何变化。
\end{example}

\begin{proof}[解]
与电源断开意味着极板总电荷 $Q$ 保持不变。插入电介质后
\[
C'=\kappa C.
\]
由 $U=Q/C$ 得
\[
U'=\frac{Q}{C'}=\frac{U}{\kappa}.
\]
由于 $E=U/d$，故
\[
E'=\frac{E}{\kappa}.
\]
储能则变为
\[
W'=\frac{Q^2}{2C'}=\frac{W}{\kappa}.
\]
因此在断电条件下，插入电介质会使电压、场强和能量都减小。
\end{proof}

\begin{remark}
若电容器始终接在恒压电源上，则情况不同：电压 $U$ 保持不变，电容增大为 $\kappa C$，电荷增大为 $\kappa Q$，总储能变为 $W'=\kappa W$。这类题目的关键是先判断\textbf{不变量}究竟是 $Q$ 还是 $U$。
\end{remark}

下面选取三类常见经典题，展示如何把“定义、求导、积分”真正用起来。

\section*{本章小结}

\begin{itemize}
  \item 电势差由线积分定义：$V_B-V_A=-\int_A^B \vb E\cdot \dd\vb l$；
  \item 静电场满足 $\vb E=-\grad V$，电场线与等势面垂直；
  \item 点电荷电势为 $V=\dfrac{1}{4\pi\varepsilon_0}\dfrac{q}{r}$，连续分布可积分求得；
  \item 电容定义为 $C=Q/U$，平行板电容器满足 $C=\varepsilon_0 S/d$；
  \item 电场能量可写为 $W=\dfrac{1}{2}CU^2$，能量密度为 $u=\dfrac{1}{2}\varepsilon_0E^2$；
  \item 电介质通过极化削弱内部电场，并使电容增大为 $\kappa$ 倍。
\end{itemize}

